% Modernised reading — fonds Alexandre Grothendieck, shelfmark 156-5, pages 1-48.
%
% This is an INTERPRETATION, not a transcription. Everything the critical
% apparatus carried moves into footnotes. In any dispute of substance the
% transcription governs: whoever cites Grothendieck must cite batch-NN.fr.tex.
%
% Pass: Opus 5.5 (claude-opus-5-5), 2026-10-10 — first pass, unchecked against
% the pages by a human. The folder was read whole: three batches, pages 1-48,
% all mathematics, dated 14 June 1986 (page 1), 15 June (page 17) and 18 June
% (page 36). Not measured against the Opus 5 reference reading of folder 115.
%
% What the folder is. Chapter V (« GF V ») of Vers une géométrie des formes,
% « Algèbre des figures », written between the first mouture of Analysis situs
% (GF IV = 156-4, whose last date is 14 June) and the second (GF VI = 156-6,
% begun 18 June). It builds a set-theoretic model of the axioms C 1-C 8 of the
% first mouture: a « figure ensembliste » is an admissible family of subsets of
% a set L, equivalently a part S of L with an Alexandrov (« ultra- ») topology.
% Sub-figure, raffinement and subdivision become closed subspace, continuous
% inclusion and finer topology; the axioms are checked one by one; the
% extension of a subdivision (C 7) becomes a gluing of topologies, whose
% characterisation is attempted, found false (pp. 42, 47, 48), and corrected.
%
% Notation fixed for the whole document. L the ambient set. F, G, H, F', G'
% figures (pages 1-10 write 𝔏, Ψ, Σ for them; pages 24-40 write 𝓛 for the
% ambient set). |F| the support, ∂A and A° (A° = A ∖ ∂A, the open stratum),
% A_F(x) the least member of F containing x, R_F the partition into open
% strata. 𝔉 the set of figures (𝒥 on page 24), 𝓜 the multiplicities. ≤
% sub-figure, ≪ raffinement, ≼ subdivision, F.G the intersection figure.
% Topologies are written τ throughout (pages 36-40 write T, 𝒯); the
% specialisation preorder is x ≤ y ⇔ x ∈ closure(y), so closed = down-closed.
% In the gluing: S = |F|, T = |G| closed, U = S ∖ T, τ'_0 on T, τ' the glued
% topology, τ'' a competitor.
%
% Substantive departures from the pages, each footnoted where it occurs:
% — p. 3: the struck counter-example is wrong (B = L contains A); his margin
%   sends to GF VIII his pp. 44-45 (= 156-8 pp. 52-53) for a correct one.
% — p. 5: for infinite families « the boolean algebra generated » must be the
%   complete one (arbitrary unions).
% — p. 11: « irreducible ⇔ greatest element » holds for complete
%   join-irreducibility in 𝔉; the topological aside fails for infinite down-sets.
% — pp. 12-14: « 𝓜 is connected » is FALSE as soon as card L ≥ 2 ({L} is
%   isolated; {L∖{x}} and {L∖{x}, L} form a component), and C 4 fails ({L},
%   {{x}} incompatible). The multiplicities of pp. 12, 13, 14 are not
%   admissible and the formula (A ∪ B)° = A° ∪ B° for ⊠ is false. True when the
%   members are finite subsets of an infinite L. Our analysis, footnoted.
% — p. 15: « disjoint » is GF IV's (compatible, no common member).
% — p. 17: the density margin is an equivalence only for closed members.
% — p. 23: the margin's d) without the closedness clause is not enough.
% — p. 24: commutativity of the square supplied.
% — p. 28: (𝔉, ≪) is even a complete lattice ({L} is its top); for an infinite
%   family the inf carries the Alexandrov topology of the meet preorder, not
%   the initial topology of p. 38's margin.
% — pp. 31-33: the « lemma » used in d ⇒ b is false in both directions; b ⇔ c
%   holds, d is independent of them; the crossed-out a) is TRUE (F' is the
%   greatest subdivision inducing G'), and F' is characterised by c) + d).
% — p. 34: the words « croissante » / « ensembliste » for g_J, h_J swapped.
% — pp. 41-48: e) is false, as he finds; counter-example of p. 48 made
%   explicit; his closing condition is right and is proved here; the closure
%   formula has ∪ (p. 44 margin), not ∩ (p. 48); p. 43's ≠ is =; p. 45's
%   « partie de T » is a part of S.
%
% Modern names given and footnoted as ours: topologie d'Alexandrov et préordre
% de spécialisation (Alexandroff 1937), espace stratifié au-dessus d'un
% ensemble ordonné (Lurie, after 1986), ouverts réguliers (Stone 1937),
% condition de frontière (Whitney, Mather), complexe CW régulier et ensemble
% de ses faces (Björner 1984), treillis complet, recollement (SGA 4, exposé IV),
% théorème A de Quillen (1973) and « foncteur asphérique ».
%
% What the folder announces and never establishes: e) of p. 41 (false); the
% equivalence (b, c, d) ⇔ a) hoped for on p. 36 (false as stated); the
% axioms C 8 (« trivial », p. 26, never stated here); the passage to
% topological spaces of p. 16 (fragmentary).
\documentclass[11pt,a4paper]{article}
\input{../preamble/grothendieck}

\folder{156-5}
\pages{1}{48}
\dating{1986}
\watermark{Édition de démonstration\\interprétation personnelle de l'œuvre}
\foldertitle{[Chapitre] V. Algèbre des figures : notes manuscrites (14/06/1986) — lecture modernisée du dossier entier}

\begin{document}

\section*{Résumé}

\begin{resume}
Une carte de géographie découpe un territoire en pays ; une triangulation
découpe une surface en sommets, arêtes et faces. Dans les deux cas on a des
morceaux, chacun bordé par des morceaux plus petits : une face est bordée par
des arêtes, une arête par deux sommets. Ces quarante-huit feuillets, écrits du
14 au 18 juin 1986, cherchent ce qui reste d'un tel découpage quand on oublie
tout de la géométrie et qu'on ne garde que des ensembles.

Le chapitre précédent (dossier 156-4, première « mouture » d'une
\emph{Analysis situs}) avait posé des axiomes pour trois manières de comparer
deux découpages, qu'il appelle des \emph{figures} : l'un peut être un morceau
de l'autre (sous-figure), l'un peut être plus fin que l'autre sur le même
territoire (subdivision), et, entre les deux, l'un peut être un découpage plus
fin d'un morceau de l'autre (raffinement). Le présent chapitre construit le
modèle le plus pauvre possible de ces axiomes : une figure n'est plus qu'une
famille de parties d'un ensemble, soumise à une condition simple --- chaque
point a un plus petit morceau qui le contient, et chaque morceau a un
« intérieur » non vide, ce qui n'est pas recouvert par les morceaux plus
petits. On vérifie les axiomes un à un.

L'idée qui arrive en route, le 18 juin, est qu'une telle famille n'est rien
d'autre qu'une topologie d'un genre très particulier, où toute réunion de
fermés est fermée : les morceaux sont les adhérences des points. Les trois
comparaisons deviennent alors des notions de topologie de première année ---
sous-espace fermé, inclusion continue, topologie plus fine --- et les preuves
laborieuses du début se raccourcissent. Le dernier problème, prolonger à toute
une figure la subdivision d'un de ses morceaux, devient celui de recoller deux
topologies le long d'un fermé.

Les dernières pages montrent un auteur qui se corrige à vue. Une
caractérisation énoncée à la page 41 est déclarée « fausse » en marge de la page
42, puis « c'est faux !! » à la page 47, « ultra faux » à la page 48, où un
contre-exemple à trois points la défait ; deux lignes plus bas vient la
condition juste. Une marge de la page 42 ajoute que c'est « une condition que
j'avais d'abord oubliée à la page 31 ». La présente lecture va un peu plus loin
dans le même sens : la connexité affirmée aux pages 12 à 14 est fausse, et la
caractérisation de la page 31 l'est aussi, sauf la condition qu'il a barrée.

Les noms à chercher : topologies d'Alexandrov et préordres, espaces stratifiés
au-dessus d'un ensemble ordonné, recollement.

\keywords{Alexandrov topology, specialization preorder, poset-stratified set, closed subspace, complete lattice, regular open set, frontier condition, regular CW complex, face poset, refinement of stratifications, gluing of topologies, recollement, Quillen's Theorem A}
\end{resume}

\section*{Le fil du dossier, et les conventions}

\begin{enumerate}
\item[1.] \emph{Familles admissibles} (pages 1 à 5) : la condition, et le
dictionnaire avec les ensembles ordonnés.
\item[2.] \emph{Sous-figures et compatibilité} (pages 6 à 11).
\item[3.] \emph{L'ensemble ordonné des figures et les axiomes C 1 à C 4}
(pages 10 à 14), dont la connexité des multiplicités, qui est fausse.
\item[4.] \emph{Supports, complément, et vers les espaces} (pages 15 et 16).
\item[5.] \emph{Sous-classes de figures} (pages 17 à 19), le 15 juin.
\item[6.] \emph{Raffinement et subdivision} (pages 19 à 26), et les axiomes
C 5 à C 8.
\item[7.] \emph{La figure intersection} (pages 27 à 29), et l'axiome C 6.
\item[8.] \emph{Prolonger une subdivision} (pages 29 à 36), l'axiome C 7.
\item[9.] \emph{Le point de vue des ultra-topologies} (pages 36 à 40), le
18 juin.
\item[10.] \emph{Recollement de topologies, et une caractérisation fausse}
(pages 41 à 48).
\end{enumerate}

Sa pagination 1 à 48 coïncide avec celle des archivistes. La page 1 porte en
angle « GF V » : chapitre V de \emph{Vers une géométrie des formes}. Le
chapitre IV (dossier 156-4), première mouture de l'\emph{Analysis situs}, porte
en dernier la date du 14 juin, celle de la page 1 d'ici ; le chapitre VI
(dossier 156-6), deuxième mouture, s'ouvre le 18 juin, date que porte aussi la
marge de la page 36 d'ici, là où le point de vue change. Le chapitre VI
s'intitule d'ailleurs, à sa première page, « Algèbre de figures ». Deux marges
d'ici (pages 1 et 3) renvoient au chapitre VIII (dossier 156-8) : elles ont été
ajoutées à une relecture postérieure d'au moins dix jours\footnote{Les renvois
sont à sa propre pagination du chapitre VIII, « pp.~37 » et « pp.~44, 45 » ;
dans le dossier 156-8 ce sont les pages 45 et 52-53 des archivistes, où il écrit
« Je vais reprendre ici la prop.~1 p.~1 de GF V, dans un contexte indiciel ».}.

\emph{Ce que le dossier vérifie.} Les axiomes de la première mouture sont
nommés C 1 à C 8 ; ce sont ceux des pages 53 à 78 du dossier 156-4 (C 1 à C 3
pages 53 à 58, C 4 page 64, C 5 à C 7 pages 69 à 72, un C 8 barré page 78). On
les retrouve ici pages 10 et 11 (C 1 à C 3), 13 (C 4), 26 (C 5 à C 8), 29 (C 6
et C 7). Le dossier les tient tous pour vrais dans le modèle ensembliste ; C 4
ne l'est pas (pages 12 à 14).

\emph{Conventions.} $L$ est l'ensemble ambiant\footnote{Les pages 24 à 40
l'écrivent $\mathcal{L}$.}. Une \emph{figure} est une famille admissible de
parties de $L$ (page 3), notée $F$, $G$, $H$\footnote{Les pages 1 à 10 notent
les familles $\mathfrak{L}$, $\Psi$, $\Sigma$ ; à partir de la page 11 elles
deviennent des figures $F$, $G$. On prend la seconde notation partout, pour ne
pas confondre $L$ et $\mathfrak{L}$.} ; ses éléments sont ses \emph{strates
fermées} (ou membres), et pour $A \in F$
\[
\partial A = \bigcup_{B \in F,\ B \subsetneq A} B, \qquad A^{\circ} = A
\smallsetminus \partial A
\]
est la \emph{strate ouverte} de $A$. $|F|$ est le support, $A_{F}(x)$ le plus
petit membre de $F$ contenant $x$, $R_{F}$ la partition de $|F|$ en strates
ouvertes. $\mathfrak{F}$ est l'ensemble des figures\footnote{La page 24 l'écrit
d'une lettre lue $\mathcal{J}$.}, $\mathcal{M} \subset \mathfrak{F}$ celui des
multiplicités. $F \leq G$ (sous-figure), $F \ll G$ (raffinement), $F
\preccurlyeq G$ (subdivision). Les topologies sont notées $\tau$\footnote{Les
pages 36 à 40 écrivent $T$ puis $\mathcal{T}$ pour une topologie, alors que $T$
y désigne aussi un ensemble ; on unifie.} ; le \emph{préordre de
spécialisation} est $x \leq y \iff x \in \overline{\{y\}}$, de sorte que les
fermés sont les parties stables vers le bas --- c'est la convention de la
page 4.

\emph{Ce que le dossier annonce et n'établit pas} : la caractérisation e) de la
page 41, qu'il déclare lui-même fausse ; l'équivalence entre les conditions de
la page 31 et la maximalité, espérée page 36 ; l'énoncé de C 8, dit
« trivial » page 26 sans être écrit ; le passage aux espaces topologiques de la
page 16, qui n'est lisible que par fragments.

\pagerange{1}{5}
\section*{Familles admissibles (pages 1 à 5)}

\subsection*{La condition}

Soit $F$ un ensemble de parties de $L$. Le paragraphe s'ouvre, le
14 juin 1986, sous le titre « Sorite ensembliste »\footnote{« Sorite » est lu
avec doute. Dans la définition de $\partial A$, la page écrit $B \subset A$ sous
la réunion ; l'inclusion doit être stricte, faute de quoi $A^{\circ}$ serait
vide.}.

\textbf{Proposition 1} (pages 1 et 2). \emph{Les conditions suivantes sont
équivalentes :}
\begin{enumerate}
\item[1°)] \emph{pour tout $x \in |F|$, l'ensemble $F_{x} = \{A \in F \mid x \in
A\}$ a un plus petit élément $A_{F}(x)$ ;}
\item[2°)] \emph{($\alpha$) pour $A, B \in F$, $A \cap B$ est réunion des
membres de $F$ qu'il contient ; ($\beta$) tout $x \in |F|$ appartient à la
strate ouverte d'un membre ;}
\item[3°)] \emph{($\gamma$) les strates ouvertes non vides forment une
partition de $|F|$ ; ($\delta$) tout $A \in F$ est saturé pour cette
partition ; ($\eta$) si $A^{\circ} \neq \emptyset$ et $A^{\circ} \subset B$,
avec $A, B \in F$, alors $A \subset B$.}
\end{enumerate}

Le point est que $x \in A^{\circ}$ signifie exactement que $A$ est un élément
\emph{minimal} de $F_{x}$\footnote{La page écrit dans ($\beta$) « $\forall x \in
\mathfrak{L}$ » ; il faut lire $x \in |\mathfrak{L}|$.} : 1°) demande un plus
petit élément, 2°) un élément minimal plus la condition ($\alpha$), qui force
le minimal à être le plus petit. La démonstration de la page 2 est complète :
si $x \in A^{\circ}$ et $x \in B$, ($\alpha$) donne $C \in F$ avec $x \in C
\subset A \cap B$, la minimalité de $A$ donne $C = A$, donc $A \subset B$ ; de
là ($\gamma$), ($\delta$), ($\eta$) ; et 3°) $\Rightarrow$ 1°) parce que, si
$x \in A^{\circ}$ et $x \in B$, ($\delta$) donne $A^{\circ} \subset B$ et
($\eta$) $A \subset B$.

Une marge propose de remplacer ($\delta$) et ($\eta$) par une seule condition :
($\delta'$) tout $A \in F$ est la réunion des $B^{\circ}$ pour $B \in F$, $B
\subset A$. C'est juste\footnote{Avec ($\gamma$) : si $x \in A^{\circ}$ et $x
\in B$, ($\delta'$) appliqué à $B$ donne $C \subset B$ avec $x \in C^{\circ}$,
et ($\gamma$) donne $C = A$, donc $A \subset B$ ; c'est 1°). La marge, de sa
main à la même encre, ajoute « voir variante \emph{in extenso}, formellement
plus forte, voir GF VIII p.~37 », renvoi à une relecture postérieure.}.

\emph{Remarques} (page 3). La condition ($\beta$) est automatique si $F$ est
finie, et plus généralement si $F$ est artinienne pour l'inclusion. Elle ne
l'est pas sans cela : une suite strictement décroissante $A_{1} \supsetneq
A_{2} \supsetneq \cdots$ d'intersection non vide satisfait ($\alpha$) mais non
($\beta$) aux points de l'intersection. Une seconde remarque, qui voulait
montrer que ($\eta$) ne découle pas d'une forme affaiblie de ($\delta'$) (avec
$B^{\circ} \subset A$ au lieu de $B \subset A$), est barrée avec son
exemple\footnote{L'exemple barré ne prouve rien : il y prend $B = \{x, y, z,
t\} = L$, qui contient $A$, et conclut « on a $A^{\circ} \subset B$ mais non
que $A \subset B$ ». La marge le dit elle-même : « Exemple idiot --- mais voir
contre-exemple correct avec $I$ infini, GF VIII pp.~44, 45 » ; ce sont les
pages 52 et 53 du dossier 156-8, qui en donnent deux, construits sur une suite
décroissante infinie. Une autre marge, « Idem si $A \in \mathfrak{L} \Rightarrow
A \neq \emptyset$ ?? », reste une question.}.

\begin{quote}
\textbf{Définition} (pages 3 et 4). $F$ est \emph{préadmissible} si elle
satisfait les conditions de la proposition 1, \emph{admissible} si de plus
$A^{\circ} \neq \emptyset$ pour tout $A \in F$.
\end{quote}

Une famille admissible ne contient pas $\emptyset$, et ses membres sont
exactement les $A_{F}(x)$ : si $x \in A^{\circ}$, $A = A_{F}(x)$.

\subsection*{Le dictionnaire avec les ensembles ordonnés}

Se donner une famille admissible $F$ revient à se donner (page 4)
\begin{enumerate}
\item[a)] un ensemble ordonné $I$ (une copie de $F$, ordonnée par inclusion),
\item[b)] une partie $S$ de $L$ (le support $|F|$),
\item[c)] une application surjective $\varphi : S \to I$ ($x \mapsto
A_{F}(x)$),
\end{enumerate}
à isomorphisme près de $I$\footnote{La précision est nôtre. Dans l'autre
sens, on pose $A_{i} = \varphi^{-1}(I_{\leq i})$ ; comme $\varphi$ est
surjective, $A_{j} \subset A_{i}$ équivaut à $j \leq i$, d'où $\partial A_{i} =
\varphi^{-1}(I_{<i})$ et $A_{i}^{\circ} = \varphi^{-1}(i) \neq \emptyset$.}.
Une note de marge, en grande partie illisible, décrit de même une famille
préadmissible : le même système, plus un ensemble de parties fermées de $I$
correspondant aux membres d'intérieur vide\footnote{Ce qui se lit de la marge
s'accorde avec l'énoncé suivant, que l'on vérifie aisément et qui est la
lecture retenue : les membres d'intérieur vide d'une famille préadmissible
sont des $\varphi^{-1}(J)$, $J$ partie de $I = F'$ stable vers le bas qui n'a
pas de plus grand élément. Une « Proposition 2 » commencée juste au-dessus est
barrée.}.

Munissons $I$ de la topologie associée à l'ordre : les fermés sont les parties
stables vers le bas --- les réunions quelconques de $I_{\leq i}$ ---, les
ouverts les parties stables vers le haut ; et $S$ de la topologie image
réciproque. Alors
\[
F = \bigl\{ \overline{\varphi^{-1}(i)} = \varphi^{-1}(I_{\leq i})
\bigr\}_{i \in I} :
\]
les membres de $F$ sont les adhérences des strates ouvertes. Les fermés de $S$
sont les réunions quelconques de membres,
\[
A_{J} = \bigcup_{i \in J} A_{i} = \coprod_{i \in J} A_{i}^{\circ} =
\varphi^{-1}(J), \qquad J \subset I \text{ stable vers le bas},
\]
les ouverts sont les $\varphi^{-1}(J)$ pour $J$ stable vers le haut --- et cette
fois, prévient la page, on ne peut pas remplacer $A_{i}^{\circ}$ par $A_{i}$
---, et les parties localement fermées sont les $\varphi^{-1}(J)$ pour $J$
\emph{convexe} : $i \leq j \leq k$, $i, k \in J \Rightarrow j \in J$.

Une topologie où toute réunion de fermés est fermée --- où chaque point a un
plus petit voisinage --- est ce qu'on appelle aujourd'hui une \emph{topologie
d'Alexandrov}, et elle est la même chose qu'un préordre, le préordre de
spécialisation\footnote{P.~Alexandroff, « Diskrete Räume », 1937. La page 36
appellera ces topologies « ultra-topologies ». Le couple $\varphi : S \to I$,
continu pour la topologie d'Alexandrov de $I$, est ce que la littérature récente
appelle un ensemble (ou un espace) \emph{stratifié au-dessus de l'ensemble
ordonné $I$} --- par exemple J.~Lurie, \emph{Higher Algebra}, appendice A.5 ---,
avec la même convention que les ouverts sont stables vers le haut. Ces
références sont nôtres ; la seconde est postérieure de trente ans.}. Ce qui
s'écrit ici page 4 en termes d'ensembles d'indices sera, le 18 juin, le point
de départ de la page 36.

\emph{Parties constructibles} (page 5). Les parties de $S$ saturées pour $R_{F}$
--- réunions quelconques de strates ouvertes --- correspondent bijectivement aux
parties quelconques de $I$ par $J \mapsto \varphi^{-1}(J)$. La page les appelle
\emph{$F$-modérées} ou \emph{$F$-constructibles}, et ajoute, pour $F$ finie,
que ce sont exactement les parties obtenues à partir des membres par réunions,
intersections et complémentaires dans $S$, les $A^{\circ}$ étant les atomes de
l'algèbre de Boole ainsi engendrée\footnote{Pour $F$ infinie, c'est vrai de
l'algèbre de Boole \emph{complète} engendrée (réunions et intersections
quelconques) ; l'algèbre engendrée par des opérations finies est en général
plus petite. La page restreint bien l'énoncé à $F$ finie, par une addition au
dessus de la ligne. Elle écrit « parties $\mathfrak{L}$-constructibles de
$\mathfrak{L}$ » : ce sont des parties de $S$.}.

\pagerange{6}{11}
\section*{Sous-figures et compatibilité (pages 6 à 11)}

\subsection*{Sous-figures}

Une \emph{sous-famille admissible} de $F$ est une partie $G \subset F$
\emph{fermée}, c'est-à-dire stable vers le bas pour l'inclusion. Pour $A \in G$
on a alors $\partial_{G} A = \partial_{F} A$, donc les strates ouvertes sont les
mêmes, $G$ est admissible, et $|G| = A_{G}$ est une partie fermée et saturée
de $|F|$\footnote{Les segments de cette démonstration sont écrits en surcharge
les uns des autres ; on les donne dans l'ordre qu'impose l'argument.}.

\textbf{Proposition} (pages 6 et 7). \emph{Soient $F$ admissible et $G$ un
ensemble de parties de $L$. Les conditions suivantes sont équivalentes :}
\begin{enumerate}
\item[a)] \emph{$G \subset F$ et $G$ est fermée dans $F$ ;}
\item[b)] \emph{$G \subset F$, $G$ est admissible, et la partition $R_{G}$ de
$|G|$ est celle induite par $R_{F}$.}
\end{enumerate}

Pour b) $\Rightarrow$ a), soient $A \in G$ et $B \in F$, $B \subset A$ : la
$R_{F}$-classe $B^{\circ}$ est contenue dans $|G|$, c'est donc une $R_{G}$-classe
$C^{\circ (G)}$, et l'on en tire $C = B$, donc $B \in G$\footnote{La page ajoute
dans a), au-dessus de la ligne, « $|\Psi|$ est saturé dans $|\mathfrak{L}|$ »,
et s'en sert dans la preuve. C'est automatique : $|G|$ est réunion de membres de
$F$, qui sont saturés par ($\delta$). Le pas « cela implique $B = C$ » est
elliptique ; il se justifie en regardant un point de $C^{\circ (F)}$ si $B
\subsetneq C$, dont la $R_{G}$-classe serait à la fois $C^{\circ(G)}$ et
$C^{\circ(F)}$, disjointe de $B^{\circ}$.}.

\emph{Scholie} (page 7). Les sous-figures de $F$ correspondent aux parties
fermées de $I$, ou de $|F|$ ; elles sont stables par réunions et intersections
quelconques.

\subsection*{Figures compatibles}

\textbf{Proposition} (pages 7 à 9, complétée page 11). \emph{Soient $F$ et $G$
deux figures. Les conditions suivantes sont équivalentes :}
\begin{enumerate}
\item[a)] \emph{$F \cup G$ est admissible, et $F$, $G$ en sont des sous-figures
;}
\item[a')] \emph{$F$ et $G$ sont deux sous-figures d'une même figure ;}
\item[b)] \emph{$F \cap G$ est fermée dans $F$ et dans $G$, et $|F| \cap |G| =
|F \cap G|$ ;}
\item[c)] \emph{$|F| \cap |G|$ est fermé dans $|F|$ et dans $|G|$, et pour une
partie $A \subset |F| \cap |G|$, on a $A \in F \iff A \in G$ ;}
\item[d)] \emph{pour $A \in F$ et $B \in G$, $A \cap B$ est fermé dans $A$ pour
$F_{A}$ et dans $B$ pour $G_{B}$, et les membres de $F$ et de $G$ contenus dans
$A \cap B$ sont les mêmes.}
\end{enumerate}
\emph{On dit alors que $F$ et $G$ sont \emph{compatibles}.}

La page démontre a') $\Leftrightarrow$ a) $\Rightarrow$ b) $\Rightarrow$ c)
$\Rightarrow$ a)\footnote{Dans b) $\Rightarrow$ c), la page écrit « comme
$\Psi$ fermé dans $\Psi$ » là où l'argument demande « $\mathfrak{L} \cap \Psi$
fermé dans $\mathfrak{L}$ », et dans l'énoncé de c) « $\Psi$-fermé dans $\Psi$ »
pour « dans $|\Psi|$ ». La fin de la page 8 est chargée de ratures ; le
raisonnement se reconstitue sans ambiguïté à la page 9. La répartition entre a)
et a'), écrits l'un dans l'autre à la page 7, est incertaine ; on donne celle
qui rend a') $\Leftrightarrow$ a) évident.}. Le cœur est c) $\Rightarrow$ a) :
pour $x \in T = |F| \cap |G|$, les plus petits membres $A_{F}(x)$ et $A_{G}(x)$
sont contenus dans le fermé $T$, donc appartiennent aux deux familles, donc
sont égaux ; et un membre de $G$ contenu dans un membre de $F$ est contenu dans
$T$, donc appartient à $F$, si bien que $F$ est fermée dans $F \cup G$.

La condition d) est ajoutée page 11 comme « remarque à la prop.\ p.~7 », sans
démonstration\footnote{Elle est juste. c) $\Rightarrow$ d) se lit dans $F \cup
G$. d) $\Rightarrow$ c) : $T$ est réunion des $A \cap B$, fermés, donc fermé
(une réunion quelconque de fermés l'est ici) ; et si $C \in F$ est contenu dans
$T$, un point $x \in C^{\circ}$ est dans un $B \in G$, le fermé $C \cap B$ de
$F_{C}$ contient $x$ donc $C$, d'où $C \in F_{C \cap B} = G_{C \cap B}$.}. Son
intérêt est d'être \emph{locale} : elle ne fait intervenir que les
multiplicités $F_{A}$ et $G_{B}$. C'est elle qui démontre l'axiome C 3
ci-dessous.

\textbf{Corollaires} (pages 9 et 10). \emph{Si $F$ et $G$ sont compatibles,}
\begin{enumerate}
\item[1)] \emph{$|F| \cap |G|$ est saturé dans $|F|$ et dans $|G|$, et les deux
partitions qu'il hérite coïncident ;}
\item[2)] \emph{$|F| \subset |G| \iff F \subset G$, et $|F| = |G| \iff F = G$ ;}
\item[3)] \emph{pour $F' \leq F$ et $G' \leq G$, $F'$ et $G'$ sont compatibles
et $F' \cup G'$ est une sous-figure de $F \cup G$.}
\end{enumerate}

La réciproque de 1) est fausse, prévient la page : $F$ et la famille $F^{\circ}
= \{A^{\circ}\}$ ont même support et même partition, mais ne sont compatibles
que si elles sont égales, c'est-à-dire si les membres de $F$ sont deux à deux
disjoints\footnote{C'est le corollaire 2. Que $F = F^{\circ}$ force les membres
à être disjoints vient de ($\alpha$). La phrase qui suit « En effet » s'arrête
sur une rature.}.

\pagerange{10}{14}
\section*{L'ensemble ordonné des figures et les axiomes C 1 à C 4 (pages 10 à 14)}

Soit $\mathfrak{F}$ l'ensemble des figures de $L$, ordonné par $F \leq G$ si $F$
est une sous-figure de $G$. La marge les appelle « figures
ensemblistes »\footnote{Le mot « figures » est souligné. Le reste de la page 10
est barré : il montrait que $\emptyset$ est la seule figure de support vide, et
qu'une figure compatible avec deux figures compatibles l'est avec leur
réunion.}. On a :

\begin{quote}
\textbf{C 1} (page 10). \emph{Toute famille majorée de $\mathfrak{F}$ a une
borne supérieure} : si les $F_{i}$ sont des sous-figures d'une même $G$, leur
réunion est une sous-figure de $G$, et c'est le sup.

\textbf{C 2} (page 11). \emph{Toute figure $F$ est la borne supérieure des
$F_{X} = \{Y \in F \mid Y \subset X\}$, $X \in F$.}

\textbf{C 3} (page 11). \emph{$F$ et $G$ sont compatibles --- c'est-à-dire
majorées --- si et seulement si $F_{X}$ et $G_{Y}$ le sont pour tous $X \in F$,
$Y \in G$.}
\end{quote}

Les $F_{X}$ sont les figures qui ont un plus grand élément ; la page les
appelle \emph{multiplicités ensemblistes}, et leur ensemble est noté
$\mathcal{M}$. Ce sont les éléments irréductibles de $\mathfrak{F}$, au sens où
une figure qui n'est pas borne supérieure de figures strictement plus petites
est une multiplicité, et réciproquement\footnote{La page justifie par une
parenthèse topologique : « un ensemble fermé est irréductible ssi il a un plus
grand élément ». C'est vrai pour un ensemble ordonné fini, non en général : dans
$\mathbf{N}$ muni de la topologie d'Alexandrov, $\mathbf{N}$ lui-même est un
fermé irréductible sans plus grand élément. L'énoncé sur $\mathfrak{F}$, lui,
est juste si « irréductible » s'entend : qui n'est la borne supérieure d'aucune
famille de figures strictement plus petites.}. C 3 résulte de la condition d)
de compatibilité, qui est locale.

\subsection*{Composantes connexes de $\mathcal{M}$}

L'axiome C 4 de la première mouture demande que deux multiplicités situées dans
des composantes connexes différentes de $\mathcal{M}$ soient compatibles. Les
pages 12 à 14 veulent montrer qu'ici $\mathcal{M}$ est connexe, ce qui rendrait
C 4 vide ; les composantes s'entendent pour l'ensemble ordonné $(\mathcal{M},
\leq)$, deux multiplicités étant reliées dès qu'elles ont un majorant
commun.

Les éléments minimaux de $\mathcal{M}$ sont les $\{A\}$, $A \neq \emptyset$, et
toute multiplicité qui a un membre minimal $A$ est dans la composante de
$\{A\}$. Jusque-là la page est juste. Mais l'énoncé principal --- \emph{si
$\operatorname{card} L > 1$, les multiplicités ayant un élément minimal sont
dans une même composante, et $\mathcal{M}$ est connexe au moins pour les
figures finies} (page 13), \emph{et même sans restriction} (page 14) --- est
faux.

\emph{Ce qui est vrai.} Pour $L \neq \emptyset$, la multiplicité $\{L\}$ est
isolée dans $\mathcal{M}$ : elle n'a pour sous-figures que $\emptyset$ et
elle-même, et une figure qui la contient comme sous-figure ne peut rien
contenir d'autre. Pour $x \in L$, si $L \smallsetminus \{x\} \neq \emptyset$, les
deux multiplicités $\{L \smallsetminus \{x\}\}$ et $\{L \smallsetminus \{x\},
L\}$ forment à elles seules une composante. Si $\operatorname{card} L \geq 3$,
toutes les autres multiplicités sont dans une même composante\footnote{Démonstration
(nôtre). Si $X$ a pour plus grand élément $G$ et que $L \smallsetminus G$
contient deux points $y \neq z$, alors $X \cup \{\{y\}, G \cup \{y, z\}\}$ est
une multiplicité, qui majore $X$ et $\{\{y\}\}$ ; et $\{\{y\}\}$, $\{\{y'\}\}$ sont
majorées par $\{\{y\}, \{y'\}, \{y, y', z\}\}$ pour un troisième point $z$. Une
multiplicité de plus grand élément $L$ ou $L \smallsetminus \{x\}$ qui n'est pas
l'une des trois exceptions a un membre $A$ avec $\operatorname{card}(L
\smallsetminus A) \geq 2$, et majore $F_{A}$. Enfin une multiplicité contenant
$L \smallsetminus \{x\}$ comme membre minimal ne peut avoir d'autre membre que
$L$ : un membre contenant $x$ et différent de $L$ serait $\{x\}$, et
l'intérieur de $L$ serait vide. Pour $\operatorname{card} L = 2$, $L = \{x,
y\}$, on trouve trois composantes : $\{\{L\}\}$, $\{\{\{x\}\}, \{\{x\},L\}\}$,
$\{\{\{y\}\}, \{\{y\},L\}\}$.}. Et C 4 est en défaut : $\{L\}$ et $\{\{x\}\}$
sont dans des composantes différentes et ne sont pas compatibles, puisque
$\{x\}$ est contenu dans $L$ sans être un membre de $\{L\}$.

\emph{Où l'argument se trompe.} Trois constructions ne donnent pas des familles
admissibles. La multiplicité $\{A, \{x\}, A \cup \{x\}\}$ de la page 12 a un
plus grand membre d'intérieur vide\footnote{Les accolades sont celles de la page,
$\{\{A\}, \{\{x\}\}, A \cup \{x\}\}$ ; on lit la famille évidente. Un point
supplémentaire $y \notin A \cup \{x\}$ la répare : $\{A, \{x\}, A \cup \{x,
y\}\}$. C'est ce qu'utilise la démonstration ci-dessus.}. La famille $\{\{x\}, L
\smallsetminus \{x\}\}$ de la page 13, censée relier $\{L\}$ aux autres, n'a pas
de plus grand élément et ne contient pas $L$. Enfin le produit de la page 14,
pour deux figures de supports disjoints,
\[
F \boxtimes G = \{A \cup B \mid A \in F \cup \{\emptyset\},\ B \in G \cup
\{\emptyset\}\},
\]
n'est pas admissible : $A \cup B$, pour $A$ et $B$ non vides, est la réunion de
$A$ et de $B$, qui sont des membres strictement plus petits, et son intérieur
est vide. Les formules qu'il écrit,
\[
(A \cup B)^{\circ} = A^{\circ} \cup B^{\circ}, \qquad \partial(A \cup B) =
(\partial A \cup B) \cup (A \cup \partial B),
\]
sont celles d'un \emph{produit} de cellules, $(A \times B)^{\circ} = A^{\circ}
\times B^{\circ}$ et $\partial(A \times B) = \partial A \times B \cup A \times
\partial B$, transcrites pour une réunion disjointe, où elles ne valent
plus\footnote{La même construction, avec $Z = \{B, B \cup \{x\}, \{x\} \mid B
\in X_{A}\}$, sert au haut de la page 14 et a le même défaut. La page note par
ailleurs $\operatorname{card} L \geq 1$ là où le cas traité est
$\operatorname{card} L > 1$, et $\{A\} \leq F$ pour $\{A\} \leq
\mathfrak{X}$.}.

\emph{Ce qui se sauve.} Les exceptions viennent toutes de membres trop gros,
égaux à $L$ ou à $L$ privé d'un point. Si $L$ est infini et qu'on se borne aux
figures finies \emph{dont les membres sont finis} --- c'est l'un des sens de
« figure finie » à la page 18 ---, il reste toujours deux points hors du plus
grand membre, et l'argument ci-dessus montre que $\mathcal{M}$ est connexe ; C 4
est alors vide, comme il le voulait. La page 13 parie d'ailleurs que les figures
finies « suffiront, j'imagine ».

\pagerange{15}{16}
\section*{Supports, complément, et vers les espaces (pages 15 et 16)}

Deux figures $F$ et $G$ sont \emph{disjointes} si et seulement si $|F| \cap |G|
= \emptyset$\footnote{« Disjointes » au sens de la première mouture (dossier
156-4, page 38) : compatibles, et sans membre commun. Avec ce sens l'énoncé est
juste, par c). Au sens plus faible « sans sous-figure non vide commune », il ne
l'est pas : $\{\{x, y\}\}$ et $\{\{y, z\}\}$ n'en ont aucune.}. Les figures
disjointes d'une famille $(G_{i})$ sont donc celles dont le support est contenu
dans $L' = L \smallsetminus \bigcup |G_{i}|$, et comme toute partie de $L$ est
un support, « la notion de complémentarité est la notion ensembliste ». Les
figures de support contenu dans $L'$ sont les figures de $L'$ :
$\mathfrak{F}(L') \subset \mathfrak{F}(L)$ est une partie saturée.

La page 16 dit qu'on retrouve ainsi $\mathfrak{P}(L)$ et son ordre à partir de
$(\mathcal{M}, \leq, \text{compatibilité})$, mais pas tout\footnote{Le reste
de la page, surchargé et taché d'encre, n'est lisible que par fragments. On y
devine la question de passer d'un ensemble $L$ à un espace topologique $X$, et
une distinction entre figures finies et infinies ; on ne reconstitue rien.}.
Une marge observe que $\mathfrak{F}$ est formé des parties fermées de
$\mathcal{M}$ dont les éléments sont deux à deux compatibles\footnote{C'est
juste : une réunion de figures deux à deux compatibles est encore admissible,
puisque les conditions de compatibilité sont locales.}.

Lisible enfin, la dernière ligne : dans un espace topologique $X$, les couples
d'ouverts $(U, V)$ dont chacun est l'extérieur de l'autre, $U =
\operatorname{int}(X \smallsetminus V)$ et $V = \operatorname{int}(X
\smallsetminus U)$, s'identifient aux ouverts $U$ tels que $U =
\operatorname{int}(\overline{U})$, et l'on a
\[
X \smallsetminus (U \cup V) = \overline{U} \cap \overline{V} = \partial U =
\partial V .
\]
Ce sont les \emph{ouverts réguliers}, qui forment une algèbre de Boole
complète\footnote{Résultat classique depuis M.~H.~Stone (1937) ; le
rapprochement est nôtre. La page écrit $\overline{\mathcal{U}} \cap
\overline{\mathcal{V}}$ avec des traits qui peuvent être des barres doubles.
L'idée qui se dessine --- un morceau d'espace et son complément, découpés
proprement --- ne revient pas dans ce dossier.}.

\pagerange{17}{19}
\section*{Sous-classes de figures (pages 17 à 19)}

Le 15 juin, la page 17 observe qu'on peut remplacer $\mathfrak{F}$ par divers
sous-ensembles stables par passage aux sous-figures et par borne supérieure.
Si $L$ est un espace, toute condition sur les membres $A$ pris comme parties de
$L$ en définit un : être fermés, compacts, fermés et connexes. De même toute
condition qui ne dépend que de la multiplicité $F_{A}$, comme une condition sur
$A^{\circ}$ ou sur $\partial A$ : $A^{\circ}$ connexe, ou $i$-connexe pour un
$0 \leq i \leq \infty$, $\partial A$ sphère topologique. Une marge ajoute la
condition de \emph{densité} $A = \overline{A^{\circ}}$, qu'elle dit équivalente
à : $A \subset B \Rightarrow A^{\circ} \subset \overline{B^{\circ}}$\footnote{Les
deux sont équivalentes pour une figure dont les membres sont fermés dans $L$,
non en général. La seconde est, à peu près, la \emph{condition de frontière}
des stratifications (H.~Whitney, J.~Mather) : une strate qui rencontre
l'adhérence d'une autre y est contenue. Rapprochement nôtre.}. Si $L$ porte une
structure « modérée », on ne retient que les familles de parties modérées.

Il y a aussi des conditions globales : être \emph{localement finie} (page 18),
condition stable par sup finis mais non par sup quelconques ; être
\emph{finie}, c'est-à-dire être un ensemble fini de parties ; être \emph{de type
fini}, contenue dans une réunion finie de multiplicités $F_{A}$. Une figure est
finie si et seulement si elle est de type fini et que ses strates sont
finies\footnote{« Strates » est le mot de la page ; on le comprend des
multiplicités $F_{A}$.}.

\emph{Réalisations cellulaires} (pages 18 et 19). Soit $L$ un ensemble ordonné,
réalisé par des cellules fermées $X_{x}$ ($x \in L$) d'un espace compact
triangulable $X$ : les intérieurs $X_{x}^{\circ}$ partitionnent $X$, et la
famille des $X_{x}$ est admissible --- c'est une figure de $X$. On se borne aux
parties \emph{fermées} (stables vers le bas) $A$ de $L$, pour que $A \mapsto
X_{A} = \bigcup_{x \in A} X_{x}$ soit une bijection sur les parties fermées de
$X$ saturées pour la partition en cellules ouvertes. Les figures de $L$ dont
les membres sont fermés s'interprètent alors comme des familles admissibles de
parties fermées de $X$ réunions de cellules : des modèles combinatoires de
« subdivisions » de $X$\footnote{C'est la situation d'un complexe CW régulier et
de l'ensemble ordonné de ses faces ; A.~Björner (1984) a caractérisé les
ensembles ordonnés qui en proviennent. Rapprochement nôtre. Plusieurs mots de ce
passage sont incertains, et l'adjectif « maximales » de la page 19 est
douteux.}.

\pagerange{19}{26}
\section*{Raffinement et subdivision (pages 19 à 26)}

\subsection*{Raffinement}

\textbf{Proposition} (pages 19 à 22). \emph{Soient $F$ et $G$ deux figures. Les
conditions suivantes sont équivalentes :}
\begin{enumerate}
\item[a)] \emph{pour tout $A \in F$, l'ensemble $G^{A} = \{B \in G \mid A
\subset B\}$ a un plus petit élément $\varphi(A)$ ; et pour $A \in F$, $B \in
G$, $A \cap B$ est réunion de membres de $F$ ;}
\item[b)] \emph{pour tout $A \in F$ il existe $B \in G$ avec $A^{\circ} \subset
B^{\circ}$ --- unique, noté $\varphi(A)$ ---, et $\varphi : F \to G$ est
croissante ;}
\item[b')] \emph{pour tout $A \in F$ il existe $B \in G$ avec $A^{\circ} \subset
B^{\circ}$, et pour $A \in F$, $B \in G$, $A^{\circ} \subset B \Rightarrow A
\subset B$.}
\end{enumerate}

Dans a), $\varphi$ est évidemment croissante, et a) entraîne $|F| \subset
|G|$\footnote{Une marge de la page 19 propose d'exiger $|F| \subset |G|$ ; c'est
superflu, puisque $G^{A}$ est non vide. Dans a), la page 20 écrit « $A \cap F$ »
pour $A \cap B$.}. a) $\Rightarrow$ b) : si $x \in A^{\circ}$ et $x \in B'
\subset \varphi(A)$, $B' \cap A$ est réunion de membres de $F$, l'un contient
$x$, c'est donc $A$, et $A \subset B'$ force $B' = \varphi(A)$ ; donc $x \in
\varphi(A)^{\circ}$. b) $\Rightarrow$ a), commencé page 20 et achevé page 21 :
$A = \bigcup_{A' \subset A} A'^{\circ}$ et chaque $A'^{\circ}$ est dans
$\varphi(A') \subset \varphi(A)$, donc $A \subset \varphi(A)$ ; si $A \subset
B'$, $A^{\circ}$ rencontre $\varphi(A)^{\circ}$ et $B'$, donc $\varphi(A)
\subset B'$ ; enfin $A \cap B$ est saturé pour $R_{F}$, et le lemme « $A^{\circ}
\subset B \Rightarrow A \subset B$ », écrit en marge de la page 20, permet d'y
remplacer chaque $A'^{\circ}$ par $A'$. Le corollaire b') vient de ce même
lemme (page 22)\footnote{Le signe que porte $A$ dans la ligne d'implication de
la page 22 est un astérisque ou un rond appuyé ; on le lit comme le rond de la
strate ouverte, qui n'y joue aucun rôle.}.

\begin{quote}
\textbf{Définition} (page 22). \emph{On dit alors que $F$ est un
\emph{raffinement} de $G$, et l'on écrit $F \ll G$.}
\end{quote}

C'est une relation d'ordre, et toute sous-figure est un raffinement. Une marge
de la page 21 en donne la forme la plus parlante : $F \ll G$ si et seulement si
$|F| \subset |G|$ et le préordre de $F$ sur $|F|$ est plus fin que celui que
$G$ induit\footnote{En termes de topologies d'Alexandrov --- ceux de la page 37
---, $F \ll G$ signifie que $|F| \subset |G|$ et que l'inclusion est continue.
C'est juste : la continuité dit que la trace d'un membre de $G$ est réunion de
membres de $F$, et que l'adhérence de $x$ pour $F$, qui est $A_{F}(x)$, est
contenue dans son adhérence pour $G$, qui est $A_{G}(x)$.}.

La première mouture définissait déjà le raffinement, comme deuxième ordre sur
les figures topologiques, et se heurtait à « un pb, qui a l'air intéressant »
de topologie des espaces stratifiés (dossier 156-4, pages 13 à 16). Dans le
modèle ensembliste ce problème disparaît, et la troisième mouture se souviendra
de la condition de fermeture de a) : « Dans GF V p.~19, 20, il fallait
supposer\ldots{} que pour $A \in F$, $B \in G$, $A \cap B$ soit ou bien vide, ou
une sous-strate de $A$ » (dossier 156-7, page 12).

\subsection*{Subdivision}

\textbf{Proposition} (pages 22 et 23). \emph{Pour $F \ll G$, les conditions
suivantes sont équivalentes :}
\begin{enumerate}
\item[a)] \emph{$F$ est maximal, pour $\leq$, parmi les raffinements de $G$ ;}
\item[b)] \emph{$|F| = |G|$ ;}
\item[c)] \emph{toute strate ouverte de $G$ est réunion de strates ouvertes de
$F$.}
\end{enumerate}
\emph{On dit alors que $F$ est une \emph{subdivision} de $G$, et l'on écrit $F
\preccurlyeq G$.}

b) $\Rightarrow$ a) parce qu'une sous-figure de même support est la figure
elle-même ; a) $\Rightarrow$ b) parce que, si $x \in |G| \smallsetminus |F|$,
$F \cup \{\{x\}\}$ est un raffinement de $G$ qui contient $F$ comme sous-figure
stricte --- on utilise que la réunion de deux raffinements compatibles en est un,
le raffinement se vérifiant sur les multiplicités. La condition c) est ajoutée
en marge ; une variante d) y est donnée aussi, qui demande que toute strate de
$G$ soit réunion de strates de $F$ et que chaque $G^{A}$ ait un plus petit
élément\footnote{Si l'on n'y suppose pas $F \ll G$, d) ne suffit pas : sur $L =
\{p, q\}$, $G = \{\{p\}, \{p,q\}\}$ et $F = \{\{q\}, \{p,q\}\}$ la satisfont,
mais $\{p,q\} \cap \{p\} = \{p\}$ n'est pas réunion de membres de $F$. Ce sont
les deux ordres opposés sur deux points. La marge est en partie illisible et
pourrait supposer $F \ll G$.}.

\subsection*{Le dictionnaire}

Se donner $F \ll G$, c'est se donner (page 24) un diagramme commutatif
\[
\begin{array}{ccccc}
S = |F| & \hookrightarrow & T = |G| & \hookrightarrow & L \\
\downarrow p & & \downarrow q & & \\
I & \xrightarrow{\ \varphi\ } & J & &
\end{array}
\]
où $p$, $q$ sont surjectives et $\varphi$ croissante, non nécessairement
surjective\footnote{La commutativité $q|_{S} = \varphi \circ p$ n'est pas écrite ;
elle est nécessaire, et c'est elle qui fait de $\varphi$ l'application $A \mapsto
\varphi(A)$. Toute telle donnée définit un raffinement, puisque l'image
réciproque par $\varphi$ d'une partie stable vers le bas l'est aussi.}. Les
images réciproques des parties fermées, ouvertes, localement fermées de $T$
sont fermées, ouvertes, localement fermées dans $S$ (pages 24 et 25). Une
subdivision est le cas $S = T$ : $q = \varphi p$ avec $\varphi$ surjective.
Même bijective, $\varphi$ n'est pas forcément un isomorphisme : deux ordres sur
un même ensemble, l'un plus fin que l'autre, donnent $F \preccurlyeq G$, $F \neq
G$, avec la même partition $R_{F} = R_{G}$ (marge de la page 25).

En langage actuel, une figure est un ensemble stratifié au-dessus d'un ensemble
ordonné, et le raffinement est un morphisme stratifié dont l'application
sous-jacente est une inclusion.

\begin{quote}
\textbf{Proposition} (page 25). \emph{$F \ll G$ si et seulement si $F$ est une
sous-figure d'une subdivision de $G$ : $F \leq G' \preccurlyeq G$.}
\end{quote}

On prend $G' = F \amalg \coprod_{x \in |G| \smallsetminus |F|} \{\{x\}\}$, qui
ajoute à $F$ les points manquants comme strates isolées\footnote{Une partie b),
qui caractérisait la subdivision comme raffinement maximal, est barrée : elle
redisait la proposition des pages 22 et 23.}. Ainsi $\ll$ et $\preccurlyeq$ se
déterminent mutuellement par $\leq$. La question inverse --- retrouver $\leq$ à
partir de $\ll$ et $\preccurlyeq$ --- reste posée avec deux points
d'interrogation (page 26) ; on a du moins
\[
F \leq G \ \text{et}\ F \preccurlyeq G \Longrightarrow F = G,
\]
que la page déduit dans le cadre axiomatique de C 5 et de l'existence, pour $F
\leq G$, $F \neq G$, d'une figure non vide disjointe de $F$ qui raffine $G$.

\subsection*{Les axiomes C 5 à C 8}

La page 26 les passe en revue. \textbf{C 5} --- si $F \preccurlyeq G$, une
figure est disjointe de $F$ si et seulement si elle l'est de $G$ --- est
trivial, les supports étant égaux. \textbf{C 6}, la subdivision induite sur une
sous-figure, est vérifié page 29 ; \textbf{C 7}, le prolongement d'une
subdivision, aux pages 29 à 36. \textbf{C 8} est dit « trivial » sans être
énoncé\footnote{Le seul C 8 de la première mouture est barré (dossier 156-4,
page 78) : si une multiplicité $X$ raffine $F$, elle raffine une multiplicité
sous-figure de $F$. Dans le modèle ensembliste c'est immédiat, avec la
multiplicité $F_{\varphi(A)}$ où $A$ est le plus grand membre de $X$.}.

\pagerange{27}{29}
\section*{La figure intersection (pages 27 à 29)}

Après un essai barré, la page 27 construit, pour deux figures $F$ et $G$, la
figure
\[
F.G = \{A \cap B \mid A \in F,\ B \in G,\ A^{\circ} \cap B^{\circ} \neq
\emptyset\},
\]
indexée par l'image $K$ de $|F| \cap |G| \to I \times J$\footnote{La page écrit
$I \times J$ là où l'on attend $I \times I'$, avec ses notations du moment. Le
passage barré qui précède décrivait déjà la partition de $|F| \cap |G|$ par les
$A^{\circ} \cap B^{\circ}$, borne supérieure des deux partitions induites.}.
Ses strates ouvertes sont les $A^{\circ} \cap B^{\circ}$ non vides.

\begin{quote}
\textbf{Proposition} (page 28). \emph{$F.G$ est la borne inférieure de $F$ et
$G$ pour $\ll$ ; plus généralement toute famille $(F_{i})$ a une borne
inférieure $\bigwedge F_{i}$ pour $\ll$, de support $\bigcap |F_{i}|$.}
\end{quote}

C'est en général différent de $F \cap G$ ; mais si $F$ et $G$ sont
compatibles, les deux coïncident. La page rapproche cette existence d'un inf
d'une propriété de « demi-treillis »\footnote{Lecture incertaine. C'est même un
treillis complet : la figure $\{L\}$ est le plus grand élément pour $\ll$, car
toute inclusion dans un espace grossier est continue, et un ensemble ordonné
où toute famille a un inf est un treillis complet. En termes de préordres,
$\bigwedge F_{i}$ est $\bigcap |F_{i}|$ muni de l'intersection des préordres ;
pour une famille infinie, c'est la topologie d'Alexandrov de ce préordre, plus
fine en général que la topologie initiale dont parle la marge de la page 38.}.
Une proposition barrée annonçait la « subdivision induite par un raffinement » ;
elle est reprise ainsi :

\textbf{Proposition} (page 29). \emph{Soient $F'$, $F$, $G$ trois figures.}
\begin{enumerate}
\item[a)] \emph{$F' \ll F \Rightarrow F'.G \ll F.G$ ;}
\item[b)] \emph{$F' \preccurlyeq F \Rightarrow F'.G \preccurlyeq F.G$ ;}
\item[c)] \emph{$F' \leq F \Rightarrow F'.G \leq F.G$, et alors $F'.G = \{C \in
F.G \mid C \subset |F'|\}$.}
\end{enumerate}

a) est la monotonie d'un inf ; b) s'en déduit, les supports étant égaux ; c) se
démontre en notant qu'un membre de $F'$ a la même strate ouverte dans $F'$ et
dans $F$, et que $A' \cap B' \subset A \cap B$ avec $A'^{\circ} \cap B'^{\circ}
\neq \emptyset$ force $A' \subset A$.

\begin{quote}
\textbf{Corollaire} (page 29). \emph{L'axiome C 6 est vrai : si $F'
\preccurlyeq F$ et $G \leq F$, il existe une unique subdivision $G'$ de $G$ qui
soit une sous-figure de $F'$, à savoir $G' = F'.G = \{A \in F' \mid A \subset
|G|\}$.}
\end{quote}

\noindent La page ne détaille pas le corollaire\footnote{On le tire de b) et de
c) appliqué à $G \leq F$ : $F.G = G$, et $|G|$, fermé pour $F$, l'est pour la
topologie plus fine de $F'$. Une marge encadrée formule en outre deux conditions
qu'elle appelle C 7${}_{0}$ : a) $P \ll Q \ll R$ et $P \preccurlyeq R$
entraînent $P \preccurlyeq Q$ et $Q \preccurlyeq R$ ; b) $P \ll Q \ll R$, $P
\leq R$ et $Q \leq R$ entraînent $P \leq Q$. Les deux sont vraies ici : a) parce
que les supports sont pris en étau, b) parce qu'un membre de $R$ contenu dans
$|Q|$ appartient à $Q$.}.

\pagerange{29}{36}
\section*{Prolonger une subdivision (pages 29 à 36)}

\subsection*{La subdivision prolongée}

Soient $G \leq F$ et $G' \preccurlyeq G$. L'axiome C 7 demande une subdivision
de $F$ qui induise $G'$ sur $G$. La page 29 prend
\[
F' = G' \cup (F \smallsetminus G),
\]
réunion disjointe : on remplace les membres de $G$ par ceux de $G'$, et l'on
garde les autres. C'est, mot pour mot, la formule de C 4 dans la première
version des axiomes de la première mouture (dossier 156-4, page 47).

Il faut vérifier que $F'$ est une figure, qu'elle subdivise $F$, et qu'elle
induit $G'$ (pages 30 et 31). On a $|F'| = |F|$. Pour $x \in |G|$, le membre
$A$ de $G'$ dont la strate ouverte contient $x$ est le plus petit membre de $F'$
qui contient $x$ : si $x \in B \in F \smallsetminus G$, la trace $B \cap |G|$
est réunion de membres de $G$, donc de $G'$, et contient $x$, donc $A$. Pour $x
\notin |G|$, c'est $A_{F}(x)$, qui n'est pas dans $G$\footnote{La page ne
vérifie pas que les strates ouvertes de $F'$ sont non vides. Elles le sont : un
membre de $F$ contenu dans $|G|$ est dans $G$, donc $\partial_{F'} A =
\partial_{G'} A$ pour $A \in G'$, et $\partial_{F'} A = \partial_{F} A$ pour $A
\in F \smallsetminus G$. Le signe sur $A$, page 30, est le même astérisque qu'à
la page 22.}. Le raffinement se vérifie membre à membre ; enfin $F'.G = G'$ par
la formule de c) ci-dessus.

\subsection*{Comment caractériser $F'$}

\textbf{Corollaire} (page 31), \emph{tel qu'énoncé.} \emph{Parmi les
subdivisions $F''$ de $F$ qui induisent $G'$ sur $G$, $F'$ est caractérisée
par chacune des propriétés équivalentes :}
\begin{enumerate}
\item[a)] \emph{$F'$ est le plus grand élément, pour $\preccurlyeq$, de cet
ensemble ;}
\item[b)] \emph{l'application $\varphi : F' \to F$ induit une bijection $F'
\smallsetminus G' \simeq F \smallsetminus G$ ;}
\item[c)] \emph{pour tout $A \in F \smallsetminus G$, $A^{\circ}$ est une strate
ouverte de $F'$ ;}
\item[d)] \emph{$F' \supset F \smallsetminus G$.}
\end{enumerate}

La page barre a) d'une grande croix, prouve d) $\Rightarrow$ b) (page 32) et b)
$\Leftrightarrow$ c) $\Leftrightarrow$ d) (page 33), et commence a)
$\Leftrightarrow$ (b, c, d) avant de le barrer.

\emph{Ce qui est vrai}\footnote{Analyse nôtre. On lit le corollaire comme une
liste de conditions sur une subdivision $F''$ de $F$ induisant $G'$, à
comparer avec « $F'' = F'$ ». Posons $U = |F| \smallsetminus |G|$, ouvert ; les
strates ouvertes de $F''$ contenues dans $U$ sont celles des membres de $F''
\smallsetminus G'$, et chacune est contenue dans une strate de $F$.} :
\begin{enumerate}
\item[--] a) est \emph{vrai}, et caractérise donc $F'$ : toute subdivision $F''$
de $F$ qui induit $G'$ est une subdivision de $F'$, parce que les membres de
$F'$ --- ceux de $G'$ et ceux de $F \smallsetminus G$ --- sont fermés pour
$F''$. C'est ce que les pages 36 à 40 retrouveront sous la forme « la moins fine
des topologies ».
\item[--] b) $\Leftrightarrow$ c) est vrai, par l'argument de la page 33 :
$F''$-strates et $F$-strates dans $U$ sont les fibres de deux surjections
$U \to I''^{*} \to I^{*}$, et b), c) disent toutes deux que la seconde est
bijective.
\item[--] d) n'est équivalente ni à b) ni à c). Sur $L = \{p, q\}$, avec $F =
\{\{p, q\}\}$ et $G = G' = \emptyset$, la subdivision $F'' = \{\{p\}, \{p,q\}\}$
satisfait d) et non b). Avec $F = \{\{p\}, \{p,q\}\}$ et $G = G' = \{\{p\}\}$, la
subdivision discrète $F'' = \{\{p\}, \{q\}\}$ satisfait b) et c), non d).
\item[--] c) et d) \emph{ensemble} caractérisent $F'$ : d) donne $F'' \supset
F'$, et c) interdit à un membre de $F \smallsetminus G$ d'avoir dans $F''$ une
strate ouverte plus petite que dans $F$.
\end{enumerate}

L'erreur est dans le « lemme » invoqué pages 32 et 33 : pour $F' \ll F$ et $A
\in F$, « $A \in F'$ si et seulement si $A^{\circ (F)}$ est une strate ouverte
de $F'$ ». Les deux exemples précédents le contredisent, l'un dans chaque
sens\footnote{Il est annoncé « plus bas » et n'est jamais écrit. La note
verticale de la page 32, « va se simplifier, cf plus [loin] », et la marge de la
page 33, « Attention, $\varphi^{*}$ bijectif, $\varphi^{*}$ pas néc.\ iso !!! »,
montrent qu'il sentait la difficulté ; c'est sans doute l'oubli dont parle la
marge de la page 42.}.

\subsection*{La même chose en termes d'ensembles d'indices}

Les pages 34 à 36 reprennent la situation avec $f : S \to I$ pour $F$, une
partie fermée $J \subset I$ pour $G$, $T = f^{-1}(J)$. Se donner $G'
\preccurlyeq G$, c'est factoriser $f_{J} = f|_{T}$ en $T \xrightarrow{g_{J}} J'
\xrightarrow{h_{J}} J$, avec $g_{J}$ surjective et $h_{J}$ croissante
surjective ; se donner $F' \preccurlyeq F$, c'est factoriser $f$ en $S
\xrightarrow{g} I' \xrightarrow{h} I$ ; et $F'$ induit $G'$ s'il existe $i' : J'
\hookrightarrow I'$ rendant le diagramme commutatif, avec $J' =
h^{-1}(J)$\footnote{Pour $g_{J}$ et $h_{J}$, la page écrit d'abord « appl.\
croiss.\ surjective d'ens.\ ordonnés » et « application ensembliste surjective »
dans l'ordre inverse ; la phrase suivante, pour $g$ et $h$, les remet dans le
bon ordre.}.

Il construit alors (page 35) l'ensemble ordonné
\[
K = J' \amalg (I \smallsetminus J),
\]
ordonné par l'ordre de $J'$, celui de $I \smallsetminus J$, et, pour $\alpha \in
J'$, $\beta \in I \smallsetminus J$, $\alpha \leq \beta \iff h_{J}(\alpha) \leq
\beta$ dans $I$. C'est l'ensemble d'indices de $F'$. L'application $h_{K} :
K \to I$, égale à $h_{J}$ sur $J'$ et à l'identité ailleurs, induit un
isomorphisme $K \smallsetminus J' \simeq I \smallsetminus J$, et $J' =
h_{K}^{-1}(J)$. La page 36 affirme que $K$ est objet final parmi les carrés
commutatifs de ce type, ce qui est juste\footnote{Pour $h'' : I'' \to I$
croissante avec $h''^{-1}(J) = J'$, l'application $I'' \to K$ égale à l'identité
sur $J'$ et à $h''$ ailleurs est croissante, et c'est la seule possible. C'est
la forme indicielle de a).}, et conclut : « Ça devrait impliquer (b, c, d)
$\Leftrightarrow$ a) ». Ce n'est pas le cas, on l'a vu : ce qu'il faut à côté de
b), ce sont les relations croisées $\alpha \leq \beta$ de $K$, et c'est
précisément ce qui sera trouvé page 48.

\pagerange{36}{40}
\section*{Le point de vue des ultra-topologies (pages 36 à 40)}

Le 18 juin, la page 36 change de langage : « le point de vue des
ultra-topologies est peut-être le plus pratique, sans prendre la peine des
types d'explicitation laborieuse en termes d'ensembles d'indices ordonnés ».
Une \emph{ultra-topologie} est une topologie où toute réunion de fermés est
fermée, autrement dit où tout point a un plus petit voisinage : une topologie
d'Alexandrov, c'est-à-dire un préordre.

Une figure est alors un couple $(S, \tau)$ formé d'une partie $S \subset L$ et
d'une ultra-topologie sur $S$ --- les membres sont les adhérences des points ---,
et tout ce qui précède se relit :
\begin{enumerate}
\item[--] $(S, \tau)$ et $(S', \tau')$ sont \emph{compatibles} si $\tau$ et $\tau'$
induisent la même topologie sur $S \cap S'$, fermé dans les deux ; il y a alors
une unique topologie sur $S \cup S' = S \amalg_{S \cap S'} S'$ qui induise les
deux et où $S$, $S'$ sont fermés : la topologie finale, « opération
topologique bien familière » (pages 36 et 37) ;
\item[--] les \emph{sous-figures} de $(S, \tau)$ sont ses parties fermées, avec la
topologie induite ;
\item[--] $(S', \tau') \ll (S, \tau)$ signifie $S' \subset S$ et l'inclusion
continue ;
\item[--] $(S', \tau') \preccurlyeq (S, \tau)$ signifie $S' = S$ et $\tau'$ plus fine
que $\tau$ (page 38) ;
\item[--] la borne inférieure $\bigwedge (S_{i}, \tau_{i})$ est $\bigcap S_{i}$ muni
du préordre intersection (marge de la page 38)\footnote{La marge dit
« topologie initiale », et précise aussitôt qu'elle « correspond à la notion de
relation de préordre induite par une famille de relations de préordres ». Pour
une famille finie les deux coïncident ; pour une famille infinie, la topologie
initiale n'est pas toujours d'Alexandrov, et c'est le préordre qui donne
l'inf.}.
\end{enumerate}
Les preuves des pages 6 à 33 deviennent, dans ce langage, des remarques de
topologie générale.

\subsection*{La subdivision prolongée comme topologie}

Revenant à $G \leq F$, $G' \preccurlyeq G$ (page 38), on pose $S = |F|$, $T =
|G|$, fermé dans $(S, \tau)$, et $\tau'_{0}$ la topologie de $G'$ sur $T$, plus
fine que celle induite par $\tau$. Les subdivisions de $F$ qui induisent $G'$
sont les topologies $\tau''$ sur $S$ plus fines que $\tau$ et induisant
$\tau'_{0}$ : la composée $(T, \tau'_{0}) \to (S, \tau'') \to (S, \tau)$ est
continue et l'on cherche le terme du milieu, « situation mixte » entre
topologies initiales et finales. Les fermés $A \cup B$, $A$ fermé pour $\tau$
et $B$ fermé de $(T, \tau'_{0})$, forment une ultra-topologie qui convient
(page 39) ; elle induit $\tau$ sur $S \smallsetminus T$\footnote{La page 39
vérifie la stabilité par intersections quelconques par une formule de
distributivité. Un paragraphe qui voulait montrer que toute $\tau''$ convenable
a cette forme est barré, sur les pages 39 et 40 ; il ne pouvait aboutir,
puisque c'est faux. La marge de la page 39 donne déjà la description des
fermés du lemme qui suit.}. D'où :

\textbf{Lemme} (page 40). \emph{Soient $(S, \tau)$ un espace topologique, $T
\subset S$ un fermé, $U = S \smallsetminus T$, et $\tau'_{0}$ une topologie sur
$T$ plus fine que la topologie induite. Soit $\mathfrak{T}$ l'ensemble des
topologies sur $S$ plus fines que $\tau$ et induisant $\tau'_{0}$ sur
$T$\footnote{La lettre rendue $\mathfrak{T}$ est un $\mathcal{T}$ bouclé et
souligné ; les pages 42 à 48 l'écrivent $\underline{\mathcal{C}}$.}.}
\begin{enumerate}
\item[(a)] \emph{$\mathfrak{T}$ a un plus petit élément $\tau'$ (la moins fine) ;}
\item[(b)] \emph{une partie $X$ de $S$ est fermée pour $\tau'$ si et seulement si
1°) $X \cap U$ est fermé dans $(U, \tau|_{U})$, 2°) $X \cap T$ est fermé dans
$(T, \tau'_{0})$, 3°) $\overline{X \cap U}^{(\tau)} \cap T \subset X \cap T$ ;}
\item[(c)] \emph{les fermés de $\tau'$ correspondent bijectivement aux couples
$(B, A_{U})$, $B$ fermé de $(T, \tau'_{0})$, $A_{U}$ fermé de $(U, \tau|_{U})$,
tels que $\overline{A_{U}}^{(\tau)} \cap T \subset B$ ;}
\item[(d)] \emph{si $\tau$ et $\tau'_{0}$ sont des ultra-topologies, $\tau'$ en
est une, de préordre : pour $x, y \in T$, $x \leq_{\tau'} y \iff x
\leq_{\tau'_{0}} y$ ; pour $y \in U$, $x \leq_{\tau'} y \iff x \leq_{\tau} y$ ;
jamais $x \leq_{\tau'} y$ pour $x \in U$, $y \in T$.}
\end{enumerate}

\noindent Les points (c) et (d) sont en haut de la page 41, où la page
s'ouvre au milieu de l'énumération\footnote{(c) s'interrompt une première fois
au bas de la page 40, barré. Dans (d), l'indice du premier cas est lu avec
doute $\tau'_{0}$ ; c'est le seul qui ait un sens.}. La démonstration
(pages 42 et 43) est complète : les conditions de (b) définissent une topologie,
plus fine que $\tau$, qui induit $\tau'_{0}$ ; et tout fermé pour $\tau'$
s'écrit $X = (X \cap T) \cup \overline{X \cap U}^{(\tau)}$, réunion d'un fermé
de $T$ --- donc de $S$ pour toute $\tau'' \in \mathfrak{T}$, $T$ étant fermé ---
et d'un fermé de $\tau$, donc fermé pour toute $\tau''$. On en tire la formule
d'adhérence, écrite en marge de la page 44 :
\[
\overline{X}^{(\tau')} = \overline{X \cap U}^{(\tau)} \cup \overline{X \cap
T}^{(\tau'_{0})} ,
\]
et en particulier, pour $A \subset U$, $\overline{A}^{(\tau')} =
\overline{A}^{(\tau)}$ (corollaire de la page 44)\footnote{Page 43, la ligne
« Comme $\tau'|_{\mathcal{U}} \neq \tau|_{\mathcal{U}}$ » porte un signe lu
$\neq$ ; l'argument demande l'égalité, que (b) donne.}.

C'est, pour des ensembles préordonnés, la situation du \emph{recollement} d'un
ouvert et d'un fermé complémentaire : une topologie sur $S$ dans laquelle $T$
est fermé est déterminée par ses traces sur $T$ et sur $U$ et par la façon dont
l'adhérence envoie les fermés de $U$ dans $T$ ; $\tau'$ est le recollement de
$\tau'_{0}$ et de $\tau|_{U}$ le long de l'application $A \mapsto
\overline{A}^{(\tau)} \cap T$\footnote{Le mot et la notion, pour les topos,
sont dans SGA 4 (exposé IV, § 9), dont Grothendieck est l'un des auteurs ; la
page ne fait pas le rapprochement, qui est nôtre. En termes d'ensembles
d'indices, $\tau'$ est l'ensemble ordonné $K$ de la page 35.}.

\pagerange{41}{48}
\section*{Recollement, et une caractérisation fausse (pages 41 à 48)}

\subsection*{L'énoncé e)}

Sous les hypothèses de (d), la page 41 passe aux ensembles ordonnés quotients :
$I$, $I'$ ceux de $\tau$, $\tau'$, $J \subset I$ la partie fermée qui correspond
à $T$, $\varphi : I' \to I$, $J' = \varphi^{-1}(J)$, $U' = I' \smallsetminus J'$.
Elle énonce :

\begin{quote}
\textbf{e)} (pages 41 et 42). \emph{La topologie de $I'$ est caractérisée par :
(i) elle est plus fine que l'image réciproque $\widetilde{\tau}$ de celle de
$I$, et induit sur $J'$ la topologie donnée ; (ii) elle coïncide avec
$\widetilde{\tau}$ sur $U'$ ; (iii) $\varphi$ induit un isomorphisme d'ensembles
ordonnés $U' \simeq I \smallsetminus J$ ; (iv) la topologie de $I$ est quotient
pour $I' \to I$ et $J \to I$.}
\end{quote}

\noindent et la marge de la page 42, écrite après coup, dit aussitôt « C'est
faux », avec un renvoi à la page 48\footnote{La marge est en grande partie
illisible : « C'est faux, sauf (iii), (iv) [\ldots] ! --- i.e.\
$\overline{\tau}' \in \mathcal{C}$ », et au-dessus « v.~48 », lu avec doute.
Dans (iii) la page écrit $\mathcal{U} = S \smallsetminus T$ là où l'on attend
$I \smallsetminus J$. Une autre marge, à la même hauteur : « C'est une condition
que j'avais d'abord oubliée à la page 31 ! ». Le paragraphe qui précède
l'énoncé, à la page 41, n'est lu que par fragments.}. La propriété (iv), dans
sa forme sur $S$ --- une partie de $S$ est fermée pour $\tau$ si et seulement
si elle l'est pour $\tau'$ et que sa trace sur $T$ l'est pour $\tau$ ---, est
vraie de $\tau'$ ; c'est le corollaire de la page 45, démontré par la formule
d'adhérence\footnote{La page écrit « une partie de $T$ » ; la suite (« sa trace
sur $T$ ») demande une partie de $S$. La démonstration est barrée de deux longs
traits, mais elle est juste.}.

\subsection*{Deux corollaires homotopiques}

La page 46 en tire des corollaires sur les ensembles ordonnés, notant ici
$\varphi^{*} = h_{J} : J' \to J$ et $\varphi = h_{K} : K \to I$. Si $\varphi^{*}$
est surjective, $\varphi$ l'est ; si les fibres de $\varphi^{*}$ sont
0-connexes, celles de $\varphi$ aussi\footnote{Le premier corollaire de la
page 46 s'arrête sur un mot illisible ; on le complète de la seule manière
évidente. « 0-connexes » est lu avec doute.}. Et surtout :

\begin{quote}
\textbf{Corollaire} (page 46). \emph{Si les $\varphi^{*-1}(J_{\leq x})$, $x \in
J$, sont faiblement contractiles, les $\varphi^{-1}(I_{\leq x})$, $x \in I$, le
sont aussi.}
\end{quote}

\noindent Pour $x \in J$ c'est l'hypothèse ; pour $x \notin J$, la page montre
que $\varphi^{-1}(I_{\leq x})$ a un plus grand élément, l'unique antécédent $x'$
de $x$, et un ensemble ordonné qui a un plus grand élément est contractile.
« OK ! » L'argument utilise les relations croisées de $K$ --- $\alpha \leq x'$
dès que $h_{J}(\alpha) \leq x$ --- et vaut donc pour la topologie $\tau'$, non
pour une topologie quelconque satisfaisant (i)--(iv). La condition sur les
$\varphi^{-1}(I_{\leq x})$ est l'hypothèse du \emph{théorème A} de Quillen, et
la page la qualifie d'« homotopiquement propre » ; d'où, ajoute-t-elle, une
équivalence d'homotopie faible universelle\footnote{D.~Quillen, « Higher
algebraic K-theory I », 1973. Grothendieck connaissait cette situation : \emph{À
la poursuite des champs} (1983) en fait la notion de foncteur asphérique. La page
ne nomme ni l'une ni l'autre ; le rapprochement est nôtre.}.

\subsection*{« C'est ultra faux »}

La page 47 s'aperçoit que e) n'est pas démontré : « Mais je n'ai pas démontré
\ldots{} pour prouver e), il faut montrer que $\tau'$ a les conditions dites »,
avec en marge « vrai ? c'est faux !! ». Elle reformule sur $S$ : si $\tau''$
vérifie (i) $\tau''$ plus fine que $\tau$, (ii) $\tau''|_{T} = \tau'_{0}$, (iii)
$\tau''|_{U} = \tau|_{U}$, (iv) $\tau$ est finale pour $(S, \tau'') \to S$ et
$(T, \tau|_{T}) \to S$, a-t-on $\tau'' = \tau'$ ? On sait $\tau''$ plus fine que
$\tau'$ ; pour la réciproque il faudrait la condition 3°), donc que
$\overline{A}^{(\tau'')} = \overline{A}^{(\tau)}$ pour $A \subset U$, et (iv) ne
le donne pas. Page 48 : « C'est ultra faux ».

Le contre-exemple est dessiné : $I'$ a trois éléments, $I$ deux, $J$ un, $J'$
deux. On le rend explicite ainsi\footnote{Le croquis montre deux ovales, l'un
au-dessus de l'autre ; on le décrit, on ne le redessine pas. La page écrit que
l'adhérence de $\{b\}$ est $\{a, b\}$ « pour $\tau'$ » ; c'est pour la topologie
concurrente, notée ici $\tau''$, puisque pour $\tau'$ elle est égale à
l'adhérence pour $\tau$.} : $S = \{a, b, c\}$, $T = \{a, c\}$, $U = \{b\}$ ;
$\tau$ est le préordre où $a$ et $c$ sont équivalents et inférieurs à $b$, de
sorte que $I = \{j < u\}$ ; $\tau'_{0}$ est l'ordre discret sur $\{a, c\}$. La
topologie recollée $\tau'$ est l'ordre $a < b$, $c < b$. La topologie $\tau''$
d'ordre $a < b$ seul, $c$ isolé, vérifie (i), (ii), (iii), et (iv) --- le
préordre engendré par celui de $\tau''$ et par $a \sim c$ est bien celui de
$\tau$ ---, mais l'adhérence de $\{b\}$ est $S$ pour $\tau$ et $\{a, b\}$ pour
$\tau''$. Donc $\tau'' \neq \tau'$, et e) est faux.

Sous un trait horizontal, marquée d'un double trait en marge, vient la
conclusion :

\begin{quote}
\emph{La condition à imposer pour caractériser $\tau'$ dans $\mathfrak{T}$ est
la condition évidente : pour $A \subset U$, l'adhérence de $A$ pour $\tau$ et
pour $\tau'$ est la même} --- ce qui entraîne déjà $\tau'|_{U} = \tau|_{U}$.
\end{quote}

C'est juste\footnote{Démonstration (nôtre) : si $\tau'' \in \mathfrak{T}$ et
$\overline{A}^{(\tau'')} = \overline{A}^{(\tau)}$ pour tout $A \subset U$, un
fermé $X$ de $\tau''$ vérifie 1°) et 3°) avec $A = X \cap U$, et 2°) parce que
$\tau''$ induit $\tau'_{0}$ ; il est donc fermé pour $\tau'$, et $\tau'' =
\tau'$. Pour des ultra-topologies il suffit de le demander pour les points $y
\in U$ : c'est exactement la relation croisée $x \leq y \iff x \leq_{\tau} y$
de (d), et, en termes de figures, les conditions c) et d) de la page 31
ensemble.}, et cela referme la question ouverte page 31. La dernière formule de
la page, l'adhérence d'une partie quelconque pour $\tau'$, est celle de la
marge de la page 44 :
\[
\overline{A}^{(\tau')} = \overline{A \cap U}^{(\tau)} \cup \overline{A \cap
T}^{(\tau'_{0})} .
\]
La page 48 écrit le signe du milieu $\cap$\footnote{La marge de la page 44
écrit $\cup$ ; seul $\cup$ est juste.}.

Le dossier s'arrête là. Le chapitre suivant, daté du même 18 juin que la
page 36, recommence l'\emph{Analysis situs} sous le titre d'« Algèbre de
figures ou Atelier » (dossier 156-6), où la relation $\preccurlyeq$ est dite
« formellement la plus délicate ».

\end{document}
